\scanpage{81}

\(I'_{j'},\quad (\unsure{f_{i'}}:\unsure{\mathcal E'_{i'}}\longrightarrow
\mathcal E_{\unsure{i=\tau i'}},\ \text{les }\alpha_{i'}\in\mathcal E)\)
sont pareils que dans le cas envisagé précédemment.

\uline{NB} Cela signifie en fait qu'en \unsure{commençant à} "boucher les trous" \deleted{de}
(de \(X'\)) correspondants aux \(i'\in I'\setminus I'_{j'}\), on remplace
\(\mathcal E'\) par le \deleted{\illegible{1}} groupe quotient de \(\mathcal E'\)
par le \unsure{s-groupe} invariant engendré par les
\(L_{i'}=\kappa_{i'}(T(\overline K))\)
\((i'\in I'\setminus I'_{j'})\), et les \unsure{\(\mathcal E_{i'}\)}
\((i'\in I'_{j'})\) par leurs images dans ledit groupe quotient, et on oublie
les \(I'\setminus I'_{j'}\), i.e. remplace \(I'\) par \(I'_{j'}\).

[Pour bien faire, il faudrait exprimer ces opérations aussi au niveau des diagrammes
de groupoïdes
\diagram{\begin{tikzcd}
\unsure{\Pi'_{D^*}} \arrow[r,dash] & \unsure{\Pi'_U} \arrow[r] & \unsure{\Pi'_e}
\end{tikzcd}}
\(\ldots\)].

\margin{\deleted{\unsure{Alors il faudrait} encore \unsure{expliciter}
\illegible{2} des pts de \illegible{1}}}

Cela signifie donc qu'en fait, on s'est ramené à la situation envisagée au début,
où \(S'=\unsure{\widehat f^{-1}|\widehat g^{-1}}(S)\)\ldots\ Donc finalement
la différence des deux pts de vue n'est pas immense, et
\deleted{\illegible{1}}\inserted{\unsure{celui}} adopté au début a l'avantage de
la simplicité plus grande (tout est relatif !)

\margin{80}

\scanpage{82}
\mspage{39}

11. Conditions sur les systèmes de \deleted{\unsure{de groupoïdes}}
\inserted{\unsure{gr.}} obtenus à partir de situations géométriques.

On en a déjà cité au passage, p.ex. que les groupes noyaux de
\(\Pi_{\unsure{D^*}}-\Pi_e\) sont abéliens (avec l'isom de Kummer \(\kappa\))
et que les images sont des sous-groupes \unsure{ouverts}, et de \(=\) pour
l'image des \deleted{\illegible{1}} \(\pi_1\) dans \(\Pi_U-\Pi_e\).
On va réexpliciter \inserted{et préciser} ceci dans le cadre de l'interprétation
en termes de
\diagram{\begin{tikzcd}
T(\overline K) \arrow[r,"\kappa_i"] & \mathcal E_i \arrow[r,"f_i"] &
\mathcal E \arrow[r,"p"] & \Gamma
\end{tikzcd}}
(cas de la "dim 1" --- on n'exclut pas le cas \(\unsure{I=\varnothing}\) ---
le cas "dim 0" étant trivial\ldots)

a) L'image \inserted{\(\Sigma\)} de \(p\) est ouverte, plus précisément existe
ss-groupe ouvert \(\unsure{\Gamma'|\Gamma^1}\subset\Gamma\) au dessus duquel
\(\mathcal E\) ait une section [et = "section admissible" en un sens qui sera
défini par la suite --- de sorte à exclure les sections "triviales" provenant
des \(\mathcal E_i\)\ldots].

b) L'image \(\Sigma_i\) de \(\mathcal E_i\) dans \(\Gamma\) est ouverte, et
l'extension \(\mathcal E_i\) de \(\Sigma_i\) par \(L_i\)
\((\simeq T(\overline K))\) est triviale. [En fait, on a une classe
\inserted{\uline{privilégiée}} de scindages, dits "algébriques", faisant un
torseur sous
\(K_i^*\,\unsure{\hookrightarrow}\,\varprojlim\unsure{(K_i^{*n})_n}\),
cf n° 9 \unsure{?} --- mais

\margin{81}

\scanpage{83}

on va pas considérer \inserted{pour l'instant} cet élément de structure
supplémentaire\ldots)

\deleted{c) Il faut}

Le \unsure{fort} des \deleted{hypothèses}\inserted{conditions} concerne
essentiellement la structure du groupe \(\pi=\operatorname{Ker}(p:\mathcal E\to\Gamma)\),
ses \inserted{classes de conjugaison de} \uline{sous-groupes}
\unsure{(ou plutôt} les hom extérieurs
\(\kappa_i:\deleted{\illegible{1}}T\ (=T(\unsure{\overline K^*}))\longrightarrow\pi\),
et la façon dont \(\Sigma\) opère extérieurement sur \(\pi\).
Ces conditions s'expriment de façon particulièrement simple lorsque les
\(\Sigma_i\subset\Gamma\) sont \(=\Gamma\)
("les pts de \(S\) \inserted{sont} rat / \(K\)") i.e.
\(\mathcal E_i\to\Gamma\) surj., a fortiori \(\Sigma=\Gamma\) i.e.
\(\mathcal E\to\Gamma\) surj. On se borne à ce cas. Le cas général s'en déduit
par extension finie du corps de base
[\unsure{auquel} \deleted{\illegible{2}}\inserted{\unsure{desquels}} pts de
\(S\) non rationnels sur \(K\)
\inserted{\unsure{i.e. lesquels} \(\Sigma_i\ne\Gamma\)}
\deleted{\unsure{se partagent}\illegible{1}}\inserted{donnent naissance}
à \(n_i\) pts, \deleted{\unsure{ég.}}
\(n_i=(\Gamma:\Sigma_i)\) \unsure{\(\}\)} --- et \(X_{K'}\)
\deleted{\illegible{1}} \unsure{scindé} en \(n=(\Gamma:\Sigma)\) comp. connexes
qui sont géom. connexes] --- mais j'ai la flemme d'expliciter comme il faudrait,
dans le contexte des groupoïdes au \deleted{\illegible{1}} lieu des groupes
profinis, l'opération d'extension des corps de base\ldots\ Il est entendu que

\margin{82}

\scanpage{84}
\mspage{40}

les conditions que je vais décrire seront invariantes par extension du corps de base.

\deleted{c) Il existe \(\unsure{\xi|\zeta}\in T\)
\((=\unsure{T(\overline K^*)})\) et des conjugués}

c) \(\forall i\)\deleted{\illegible{1}}, l'hom extérieur
\[
\kappa_i:T\longrightarrow\pi
\]
est compatible avec l'action extérieure de \(\Gamma\) (opérant sur \(T\) par le
caractère cyclotomique \(\chi\), et sur \(\pi\) grâce à l'extension \(\mathcal E\)
de \(\Gamma\) par \(\pi\)). En d'autres termes, pour \(\forall g\in\mathcal E\),
existe un \uline{\(\alpha\in\pi\)} (\uline{NB} pas seulement
\(\alpha\in\mathcal E\) !) tel qu'on ait,
\[
\operatorname{int}(g)\kappa_i(\unsure{\zeta})
=\operatorname{int}(\alpha)\kappa_i(\chi(\unsure{p(g)})\unsure{\zeta})
\]
i.e.
\[
\operatorname{int}(\alpha^{-1}g)\kappa_i(\unsure{\zeta})
=\kappa_i(\chi(p(g))\unsure{\zeta})
\]
Ceci signifie donc que 1°) \(\alpha^{-1}g\) normalise \(L_i=\kappa_i(T)\).
[et cela signifie \unsure{\(=\)}, probablement, que \(\alpha^{-1}g\in\mathcal E_i\)
--- et qu'on puisse trouver un tel \(\alpha\in\pi\)
\inserted{tel que \(\alpha^{-1}g\in\mathcal E_i\)} provient de l'hypothèse
\(\mathcal E_i\to\Gamma\) surjectif --- on prend un \(\beta\in\mathcal E_i\)
\inserted{\unsure{\(\alpha^{-1}g\)}} ayant \(=\) image que \(g\) dans \(\Gamma\)
et on prend \(\alpha=g\beta^{-1}\)] et que 2°) l'action \inserted{intérieure} de
\(\unsure{\beta=\alpha^{-1}g}\) sur \(T\) n'est autre que par multiplication
par \(\chi(g)=\chi(\beta)\) --- ce qui (pour \(\beta\in\mathcal E_i\)) n'est autre que la

\margin{83}

\scanpage{85}

condition déjà explicitée que l'hom de groupes \(\kappa_i:T\to\mathcal E_i\)
est compatible avec l'action de \(\mathcal E_i\), opérant sur \(T\) via
\(\chi\circ p\mid\mathcal E_i\), et sur lui-= par autom. int.

Donc la condition c) n'est pas vraiment nouvelle --- je la réexplicite en termes
un peu différents, à cause de son importance. Elle implique que l'opération
\margin{\figure{[Un signe illisible biffé est entouré d'un cercle.]}}
de \(\Gamma\) sur \(\pi\) est \uline{très} non triviale (puisque
\(\chi:\Gamma\to\widehat{\mathbb Z}^{*}\) a une image ouverte !)
--- il \unsure{n'était} pas \(=\) évident a priori
\inserted{(sans raisons arithmétiques profondes !)} --- compte tenu de la structure
de \(\pi\) qu'on va donner --- qu'il existe de telles opérations de \(\Gamma\)
sur un \(\pi\) !

Cette condition \unsure{sera} complétée par une condition de non-trivialité à la Weil.

d) \(\exists\eta\in T^*\)\unsure{.} (une base de \(T\)), et des
\(\alpha_i\in\pi\) (afin de conjuguer \(\kappa_i\) en
\(\kappa'_i=\operatorname{int}(\alpha_i)\circ\kappa_i\))
\unsure{enfin} un entier \(g\geq0\) et des éléments
\(\unsure{x_i,y_i}\in\pi\) \((1\leq i\leq g)\), tels que l'on ait

\margin{84}

\scanpage{86}
\mspage{41}

1°) les \(\kappa'_i(\eta)\), \deleted{\illegible{1}} et les
\(\unsure{x_i,y_i|x_j,y_j}\) engendrent le groupe profini \(\pi\)

2°) Ils satisfont la relation
\[
\boxed{[x_1,y_1].[x_2,y_2]\ldots[x_g,y_g]\,
\kappa'_1(\eta)\kappa'_2(\eta)\ldots\kappa'_{\unsure{\nu|n}}(\eta)=1}
\]

3°) Cette relation, avec les générateurs envisagés décrit \(\pi\)
\inserted{(en tant que groupe profini)} par générateurs et relations\ldots

\uline{NB} On sait que par ces conditions, \(g\) est uniquement déterminé,
p.ex. par le fait que
\(\pi_{ab}/\deleted{\text{images des}}\sum\unsure{\kappa_i|\delta_i}(T)\)
est un \(\widehat{\mathbb Z}\)-module libre de rang \(2g\) --- \(\pi_{ab}\)
étant libre de rang \(2g+\nu-1\) où \(\nu=\operatorname{Card}(I)\).

Pour le choix de \(\eta\), on sait que si \(\eta\) convient, alors
\(\forall\ \chi(\alpha)\eta\) aussi (où \(\alpha\in\Gamma\)) --- quitte à prendre
des conjugués. Donc les \(\eta\) qui conviennent \deleted{forment} constituent
un \unsure{sous-torseur}\inserted{de \(T^*\)} (sous le ss-groupe ouvert
\(\chi(\Gamma)\) de \(\widehat{\mathbb Z}^{*}\).

\deleted{Supposons pour simplifier que \(g=0\) :
On voit donc que si le système des \(\kappa'_i(\eta)\) est un syst. de gén.
privilégié (au sens de d) ci-dessus) alors les
\(\operatorname{int}(g)\kappa'_i(\eta)\) aussi pour \(\forall g\in\mathcal E\),
[\unsure{\(=\)} \(\operatorname{int}(\alpha)\kappa'_i(\chi(g)\eta)\)
pour \(\alpha\in\pi\) convenable]}

\margin{85}

\scanpage{87}
\mspage{42}

\margin{(\unsure{car} \(\chi:\Gamma\to\widehat{\mathbb Z}^{*}\) \unsure{épi.})}

\deleted{donc aussi les \(\kappa'_i(\chi(g)\eta)\). \unsure{Prenant}
\(K=\mathbb Q\), on trouve donc par voie arithmétique :

\uline{Prop} Soit \(\pi\) le \unsure{\(P\)}-groupe libre
\inserted{de \(\nu-1\)} \unsure{(}\uline{engendré} par des générateurs
\(\ell_1,\ldots,\ell_\nu\) satisfaisant la relation génératrice
\(\prod\ell_i=1\). Alors pour tout \(\alpha\in\widehat{\mathbb Z}^{*}\),
on a encore \(\prod\ell_i^\alpha=1\), et les \(\ell'_i=\ell_i^\alpha\)
avec cette relation décrivent le groupe profini \(\pi\).}
En fait, comme la structure du groupe \(\pi\) est indépendante de \(K\),
prenant \(K=\mathbb Q\) (et admettant qu'il existe une courbe lisse projective
géom. connexe de genre \(g\) sur \(\mathbb Q\), ayant \(\nu\) pts rat. / \(\mathbb Q\) !),
on trouve que si les \(\ell_i\), \deleted{\(\unsure{x_j,y_j}\)}
\((1\leq i\leq\nu\deleted{,\ \unsure{1\leq j\leq g}})\)
\deleted{conviennent, alors les} s'insèrent dans un syst. de générateurs
privilégié (avec des \(\unsure{x_\alpha,y_\alpha}\)) alors pour
\(\forall\rho\in\widehat{\mathbb Z}^{*}\), on peut trouver des conjugués
\(\ell'_i\) des \(\ell_i^\rho\) qui s'insèrent de \(=\).
Même pour \(\rho=-1\) ce n'est pas entièrement trivial\ldots

\deleted{\illegible{2}} Enfin, \unsure{et ce sera} une condition
\unsure{draconienne} de non trivialité de l'opération de \(\Gamma\) sur \(\pi\).
Soit \(\mathcal E'\) un sous-groupe \inserted{\uline{ouvert}}
(d'indice fini) de \(\mathcal E\),
\inserted{\(\pi'=\pi\cap\mathcal E'=\operatorname{Ker}(\mathcal E'\to\Gamma')\)}
\deleted{\illegible{1} son image dans \unsure{\(\mathcal E\)|\(\Gamma\)}
(qui est ouvert dans \(\Sigma\)) de sorte que}

\margin{86}

\scanpage{88}

\(\Gamma'\) l'image de \(\mathcal E'\) dans \(\Gamma\). On trouve donc une
extension de \(\Gamma'\) \((=\operatorname{Gal}(\overline K'/K'))\) par \(\pi'\),
donc aussi par \(\unsure{(\pi'_{ab})_{(\ell)}}\)
(\(\ell\) étant un nb premier), qui est (on le sait par d)) un
\(\mathbb Z_\ell\)-module libre de t.f.,
\inserted{sur lequel \(\Gamma'\) opère}

[Avec un peu de travail, on doit pouvoir \unsure{retrouver} sur
\(\mathcal E'\to\Gamma\) une "structure à lacets" avec des
\(\mathcal E'_{i'}\), comme pour \(\mathcal E\), et décrire dans
\(\unsure{(\pi'_{ab})_{(\ell)}}\)
\inserted{la somme} des images des \deleted{\illegible{1}} \(L'_{i'}\),
sur lesquels \(\Gamma'\) opère \inserted{donc} via le caractère \(\chi\).

\margin{Ça se fait très élégamment dans le contexte \unsure{des}
\diagram{\begin{tikzcd}
\unsure{\Pi_{D^*}} \arrow[r] & \Pi_U \arrow[r] & \Pi_e
\end{tikzcd}}}

On s'intéresse au quotient de \(\unsure{(\pi'_{ab})_{(\ell)}}\) par ces
sous-modules relativement \unsure{trivial|triviaux} (la partie "VA" du
module \(\ell\)-adique envisagé). Ceci posé

\textcircled{e} on exige que la représentation de \(\Gamma'\) ci dessus soit
"pure de poids 1" --- \deleted{\unsure{i.e.} on ait} et que les polynômes car.
des frobénius (qui sont à coefficients \(\in\mathbb Z\), pas seulement
\(\in\mathbb Z_\ell\)) soient "\uuline{indépendants}" de \(\ell\).

\margin{\figure{[Un trait vertical borde à gauche l'énoncé marqué d'un e entouré d'un cercle.]}}
\margin{87}

\scanpage{89}
\mspage{43}

12. \uline{L'analogie topologique}

Soit \(X\) une \deleted{variété}\inserted{surface} (topologique) compacte
orientable, \(S\) une partie finie de \(X\). Si \(X\) est connexe, on considère
l'ensemble \inserted{\unsure{\(\omega\)}} \unsure{des|ses} deux orientations,
on l'utilise pour tordre le groupe \(\mathbb Z\), d'où un groupe \unsure{\(T\)}
\[
T=\mathbb Z\mathop{\wedge}\limits_{\pm1}\unsure{\{\omega\}|\omega}.
\]
isomorphe (non canoniquement) à \(\mathbb Z\). Plus généralement, supposons
donné un tel groupe \(T\),
\inserted{i.e. un \(\omega\in\unsure{\operatorname{Ob}}\mathrm{Ens}_2\)}
--- une \(T\)-orientation de \(X\) est par définition un \unsure{élém.}
de l'ens. \(\operatorname{Or}(X)\mathop{\wedge}\limits_{\{\pm1\}}\omega\).

[dans le cas précédent, \(X\) sera donc canoniquement \(T\)-orienté\ldots]

Considérons le groupoïde fondamental de \(X\setminus S=U\), soit \(\Pi_U\),
et le groupoïde fondamental \inserted{\(\Pi_{D^*}\)} des germes
\unsure{disques} de \(X\) autour de \(S\), \uline{privés} de \(S\)
\inserted{\deleted{\unsure{des}} groupoïdes \unsure{?}}
(des germes de rev. universels de \(X\setminus S\) au voisinage de \(S\)).
On a donc un foncteur canonique
\diagram{\begin{tikzcd}
\unsure{\Pi_{D^*}} \arrow[r,dash] & \Pi_U
\end{tikzcd}}
et d'autre part, \deleted{\illegible{2}} si \(X\) est \(T\)-orienté, on a une
structure supplémentaire

\margin{88}

\scanpage{90}

correspondante sur \inserted{le groupoïde} \(\Pi_{D^*}\) : le système local de
ses \(\pi_1\) est can. isom. à \(T\).

Associant à tout \(X\) \(T\)-orienté le système
\[
\Pi_{D^*}\longrightarrow\Pi_{\unsure{X|U}},\qquad
\kappa:T_{(\Pi_{D^*})}\simeq\pi_1(\Pi_{D^*}/e),
\]
on trouve une 2-équivalence entre la 2 catégorie topologique des
\deleted{\unsure{variétés}} couples \((X,S)\) d'une
\inserted{(pour les homéo. à isotopie près\ldots)}
\unsure{variété compacte \(T\)-orientée}, et de la 2-catégorie des systèmes
précédents \inserted{(pour les équivalences)}, \unsure{(}qui satisfont les conditions

a) \(\pi_0(\Pi_{D^*})\), \deleted{\unsure{et}} \(\pi_0(\Pi_X)\) finis

b) Pour \unsure{tte|\(\forall\)} composante connexe de \(\Pi_U\), soit
\(\Pi_{U_0}\), et la partie \(\unsure{\Pi_{D_0^*}}\) au dessus, explicitant la
situation \unsure{groupoïdes} par \unsure{un|une}
\inserted{groupe \(\pi\) et une} famille d'homs
\(T\xrightarrow{\kappa_i}\pi\), on a (\uline{NB} le tout dépend de choix,
mais \(\pi\) est intrinsèque comme groupe extérieur, et les \(\kappa_i\)
\unsure{comme homs} extérieurs --- définissent une "\(T\)-structure à lacets
sur \(\unsure{S}\) ---) on a ce qui suit.

\margin{89}
